\section{AC-DFA Based Pattern Selection}
We first give an intuitive idea. Given a collection of
regex patterns $\{r_i\}_{i=1}^m$, the prover extracts
critical patterns for each $r_i$, and build an Aho-Corasick
automaton (AC-DFA) for the critical patterns. Each critical
pattern is associated with one or more signatures.
Given an input string $s$, the prover runs $s$ through
the AC-DFA, and get a Pedersen vector commitment to
all the final states (that capture a critical pattern).
The set of states is arranged as a concatenation of
final and non-final states, thus to verify if a state
is in final is simply a range proof.

Given a fixed set of words $\{w_i \in \Sigma^*\}$,
an Aho-Corasick automaton can be built in $O(\sum_{i=1^m} |w_i|)$ time
with $O(\sum_{i=1^m} |w_i|)$ size. An AC-DFA is deterministic
finite automaton $(S, s_0, F, T)$ where each transition $t\in T$ is a tuple
$(s_1, c, s_2, b)$ where $b$ indicates whether it is a failure edge.
Given any string $s \in \Sigma^*$, it can be accepted in 
$2|s|$ time. There exists a map $\rho$ which maps each final
state to a collection of words that the AC-DFA outputs at that state.
An AC-DFA can be further compiled into a standard
DFA with back-edges pre-computed away. 


\begin{definition} Given a regex $r$, we call a string $s$
its critical pattern, written as $s \in \critical(r)$ if
for any string $t \in L(r)$, $s$ is a substring of $t$.
\end{definition}

In another word, if a critical pattern does not appear in a string
$t$ then $t$ cannot be a match for $r$.

In practice, for $38k$ PM-Reg ClamAV rules, we collected 
$38k$ critical patterns which results in a DFA with
$2$ million states; and all patterns results in $6 million$
states for $112k$ patterns, each state has $16$ outgoing transitions
for hex alphabet (4-bit nibbles).

We will split the proof into two steps: (1) we prove a 
Pedersen commitment commits to the set set of states along
an acceptance path for an input string; (2) 
We show another Pedersen commitment commits to
the set of signatures that should be verified.

\subsection{DFA Acceptance Path}
Let $u = \log(|\Sigma|)$ and $v$ be the bit-width of
states. Given a transition $(s, c, t)$ 
 a transition can be digitized to a field element:
$\rho( (s, c, t) ) = s + 2^v c + 2^{v+u} t$.
Given the set of transitions of an AC-DFA, the
prover first pre-process $\myvec{\tau}$ as look-up table
and generate its prover keys and cached quotient information.
The set of states is encoded as a vector of consecutive numbers
which is a concatenation of final and non-final states.
Thus, to check if a state is final is translated into
a zk-range proof.

Given an input string $s$, its
acceptance path can be
modeled by four vectors: $\myvec{s}$,
$\myvec{c}$, $\myvec{t}$.
Define $\myvec{\tau}$ for each $i$ s.t.
$\myvec{\tau}_i = \rho(\myvec{s}_i, \myvec{c}_i, \myvec{d}_i)$.

Given the vector commitments to these five vectors, one can
use range-proofs to prove their relation:
\begin{enumerate}
  \item Prove $\Commit_{\myvec{\tau}} = \Commit_{\myvec{s}} + 
  2^v \Commit_{\myvec{c}} + 2^{v+u} \Commit_{\myvec{t}}$.
  \item Prove each element in $\Commit_{\myvec{s}}$,
  $\Commit_{\myvec{t}}$, $\Commit_{\myvec{d}}$ are in range,
  using batched zk-range proof.
  \item Prove that $\myvec{s}_0$ is the initial state.
\end{enumerate}

Note that given that this is a standard DFA,
$\myvec{c}$ is actually the input string $\myvec{s}$, as backedges
are computed away. Let $\vec{F}$ be the final states appears 
in the acceptance path, we then need to show that
another Pedersen commitment $\Commit_{\vec{F}}$ is valid.
This is done through the general strategy of CPL-SNARK
via fingerprinting techque. We first define some gadgets below.

\begin{enumerate}
  \item $\ttt{permutation}(\vec{a}, \vec{b}, r)$.  This is achieved
  by fingerprinting test $\prod_{i=1}^n (\vec{a}_i-r) = \prod_{i=1}^n (\vec{b}_i-r)$. Note that $\vec{a}$ and $\vec{b}$ needs to be committed first before
  random challenge $r$ is generated to apply Fiat-Shamir.
  \item $\ttt{pack}(\vec{a}, \vec{b}, \vec{c}, r)$. Given $\vec{b}$ is a boolean
  array, $\myvec{c}$ is a packed representation of $\vec{a}$, padded by
  $0$. For instance $\vec{a} = (100, 200, 300, 400)$, and
  $\vec{b} = (1, 0, 1, 0)$, and $\vec{c} = (100, 300)$. This is also
  accomplished by introducing an array $r_i$ s.t.
  $r_i = r_{i-1} * r$ iff $b_i= 1$. Then verify
  $\prod_{i=1}^n \vec{a}_i r_i$ = $\prod_{i=1}^n \vec{c}_i r^i$.
  \item $\ttt{set}(\vec{a}, \vec{b}, r)$, to test $\vec{b}$ is a 
  standard set of $\vec{a}$. The idea is to first produce
  a sorted list of $\vec{a}$, letting it be $\vec{c}$.
  Then define a boolean vector $\vec{d}$ s.t.
  $\vec{d}_i = \vec{c}_i==\vec{c}_{i-1}$. Then
  run $\ttt{pack}$ to remove the duplicates in $\vec{c}$, 
  then $\vec{b}$ is a sorted standard set of $\vec{a}$ with
  the tail padded with $0$.
\end{enumerate}

In summary, we have the construction
$\prfacc(\Commit_{\myvec{s}}, \Commit_{\myvec{f}})$ which
asserts that for string $s$, $\Commit_{\myvec{f}}$
is a vector commitment of the standard set of all
final states along its unique acceptance path over the 
AC-DFA.

{\bf Design Idea:} since we are using fingerprint checking, we
need verifier challenge. So the prover system needs to run 
in two stages (e.g., Halo2 already natively supports it).
Some complexities will arise when we need to fold relations,
the verifier challenge needs to be generated once and moved out
of IVC prover. We achieve this by aggregating/folding $k$
relations at a time, and let verifier challenge be
hash of all first stage witness wires for all instances.

Each of the relation could be modeled as a Halo2 gadget or chip
(or rely on Halo2 to optimize the use of selectors), but 
maybe it's good to have a global cusotmized layouter to
layout these relations to a standard Plonkish model (which
is later easier to do ProtoGalaxy translation). We discuss the
translation below:

$\ttt{permutation}(\vec{a}, \vec{b}, r)$ relation can be encoded
with one custom gate with $3$ columns, where an extra column
serves as the grand product. The $r$ ideally should be referenced
as a fixed cell, but it looks not possible in Halo2, and may need
an extra column. There might be an extra selector which enforces
the gate. Given there are so many instances, we may create our
own numbered selector gate to enforce it.

$\ttt{pack}(\vec{a}, \vec{b}, \vec{c})$ can be implemented similarly.
$\ttt{set}(\vec{a},\vec{b})$ is done in a similar fasion and
it needs support of range proof.

Eventually, the finite state automata acceptance logic
can be implemented by ``invoking" the above logic. However,
to save the number of rows in the constraint system, it might
be most economical to just encode all constraints together
instead of doing multiple chips and copy over the values.

\subsection{Selection from Signatures}
Now we have the set of accepstating states for
critical patterns, and then we need to link them to
the set of signatures that need to be inspected.

The trusted set up prepare a public table which 
maps from each acceptating state
$f_i$ to a vector of signature IDs.
However these signature IDs needs to be
represented compactly and we can represent them
using bilinear pairing of $[\prod_{j=1}^k(x-s_i)]_1$
for some trap-door $x$ where $\{s_{i,j}\}$  are the
related signatures to $f_i$. This table of information
is publically available.

Now the set of related signatures $\myvec{SIG}$ can
be similarlarly represented as a bilinear pairing
commitment $\Commit_{\myvec{SIG}}$.  To prove that
for state $f_i$'s signature is a subset
of $\myvec{SIG}$, one can compute the quotient polynomial
between the two characteristic polynomials and supply it
over $G_2$.

Then, we can use TIPP to batch prove this relation:
given a list of $f_i$ and the corresponding proof $\prf_i$.
All of them are hiding in TIPP's commitment to vector
of group elements. Then we use argument for inner pairing
to prove the pairing relation.

Thus, we have the construction for proving
$\Commit_{\myvec{SIG}}$ represents a {\em superset}
of the signatures related to $\Commit_{\myvec{f}}$.
Note that the prover may provide a super-set, but
it only increases the cost of the prover, and not hurting
the soundness of the proof.

\subsection{Selection from Signatures (Circuit Version)}
We have a set of final states of identifying patterns (average
size 6.6k). We have 38k signatures, and the average number of
pattens is up to 20. So we are running concurrently
38k lookups (mostly negative) and up to 50 positive.
For negative lookups, we perform $2x$ lookups (using
the trick of $a<x<b$).

In this case our lookup table is 6.6k, and our query table size
is 38k $\times 20$ = 760k. Instead of using CQ, we should
use plookup based (and range proof).

Again, given the boolean array indicates selected or not,
we can get a small packed factor which indicates which
PM-REG signature should be ruled out.


